%!TEX program = xelatex
\def\numeromodulo{1}
\def\numerolista{2}
\documentclass[12pt]{exam}

\def\ano{2023}
\def\semestre{1}
\def\disciplina{Introdução à \'Algebra Linear}
\def\turma{3}

\usepackage{caption}
\usepackage{amssymb}
\usepackage{amsmath,amsfonts,amsthm,amstext}
\usepackage[utf8]{inputenc} % Required for inputting international characters
\usepackage[T1]{fontenc}
\usepackage[brazil]{babel}
\usepackage{graphicx}
\graphicspath{{../figuras_geogebra/}{../Pictures/}}
\usepackage{enumitem}
\usepackage{multicol}
\usepackage{answers}
\usepackage[svgnames]{xcolor}
\usepackage{tikz}
\usepackage{ifthen}
\usetikzlibrary{lindenmayersystems}
\usetikzlibrary[shadings]
\usepackage{tikz-3dplot}
\usepackage{tcolorbox}
%-------------amatrix
% Augmented matrix.  Usage (note the argument does not count the aug col):
% \begin{amatrix}{2}
%   1  2  3 \\  4  5  6
% \end{amatrix}
\newenvironment{amatrix}[1]{%
  \left(\begin{array}{@{}*{#1}{r}|r@{}}
}{%
  \end{array}\right)
}

\Newassociation{solucao}{Solution}{ans}
\newtheorem{exercicio}{}

\setlength{\topmargin}{-1.0in}
\setlength{\oddsidemargin}{0in}
\setlength{\textheight}{10.1in}
\setlength{\textwidth}{6.5in}
\setlength{\baselineskip}{12mm}

\newcounter{exercicios}
\setcounter{exercicios}{0}
\newcommand{\questao}{
\addtocounter{exercicios}{1}
\noindent{\bf Exerc{\'\i}cio \Roman{exercicios}: }}

\newcommand{\resp}[1]{
\noindent{\bf Exerc{\'\i}cio #1: }}

\extraheadheight{0.7in}
\firstpageheadrule
\runningheadrule
\lhead{
        \begin{minipage}[c]{1.7cm}
        \includegraphics[width=1.7cm]{unb.pdf}
        \end{minipage}%
        \hspace{0pt}
        \begin{minipage}[c]{4in}
          {Universidade de Bras{\'\i}lia} --
          {Departamento de Matem{\'a}tica}
\end{minipage}
\vspace*{-0.8cm}
}
% \chead{Universidade de Bras{\'\i}lia - Departamento de Matem{\'a}tica}
% \rhead{}
% \vspace*{-2cm}

\extrafootheight{.5in}
\footrule
\lfoot{\disciplina\ - \semestre$^o$/\ano\ - Módulo \numeromodulo}
\cfoot{}
\rfoot{P\'agina \thepage\ de \numpages}
\newcommand{\nsub}{\varsubsetneq}
\newcommand{\vaz}{\emptyset}
\newcommand{\im}{{\rm Im\,}}
\newcommand{\sub}{\subseteq}
\renewcommand{\sin}{{\rm sen\,}}
\newcommand{\n}{\mathbb{N}}
\newcommand{\z}{\mathbb{Z}}
\newcommand{\rac}{\mathbb{Q}}
\newcommand{\real}{\mathbb{R}}
\newcommand{\complex}{\mathbb{C}}
\newcommand{\cp}[1]{\mathbb{#1}}
\newcommand{\ch}{\mbox{\textrm{car\,}}\nobreak}
\newcommand{\vesp}[1]{\vspace{ #1  cm}}
\newcommand{\compcent}[1]{\vcenter{\hbox{$#1\circ$}}}
\newcommand{\comp}{\mathbin{\mathchoice
{\compcent\scriptstyle}{\compcent\scriptstyle}
{\compcent\scriptscriptstyle}{\compcent\scriptscriptstyle}}}

\newcommand{\integer}{\mathbb{Z}}
\begin{document}

  \Opensolutionfile{ans}[ans1]
  \begin{center}
    {\Large\bf \disciplina\ - Turma \turma\ -- \semestre$^{o}$/\ano} \\ \vspace{9pt} {\large\bf
        $\numerolista^a$ Lista de Exercícios -- Módulo \numeromodulo}\\ \vspace{9pt} Prof. José Antônio O. Freitas
  \end{center}
  \hrule

\vesp{.6}

\begin{exercicio}
  Nos itens a seguir determine quais equações são lineares nas variáveis $x$, $y$ e $z$. Se alguma não for linear, explique o motivo.
  \begin{multicols}{2}
  \begin{enumerate}[label={\alph*})]
    \item $x - \pi y + \sqrt[3]{5}z = 0$
    \item $x^{-1} + 7y + z = \sin(\dfrac{\pi}{9})$
    \item $x = -7y + 3z$
    \item $3\cos(x) = 4y + z = \sqrt{3}t$
    \item $\pi x - \sqrt{x}y + \dfrac{1}{3}z = 7^{1/3}$
    \item $x^{-2} + y^2 + 8x = 5$
    \item $2x - xy - 5z = 0$
    \item $cos(3)x - 4y + z = \sqrt{3}$
  \end{enumerate}
\end{multicols}
  \begin{solucao}
    \begin{enumerate}[label={\alph*})]
      \item Sim
      \item Não
      \item Sim
      \item Não
      \item Sim
      \item Não
      \item Não
      \item Sim
    \end{enumerate}
  \end{solucao}
\end{exercicio}

\begin{exercicio}
  Nos itens a seguir, determine se as equações formam um sistema linear. Para aqueles que forem um sistema linear encontre sua solução.
  \begin{enumerate}[label={\alph*})]
    \item $\begin{cases} -2x + 4y + z = 2\\ 3x - \dfrac{2}{y} = 0\end{cases}$
    \item $\begin{cases} 4x - 4y + 2z = 0\\ -x + \ln(2)y - 3z = 0\end{cases}$
    \item $\begin{cases} x = 4\\ 2x = 8\end{cases}$
    \item $\begin{cases} -2x_1 - x_4 = 5\\ -x_1 + 5x_2 + 3x_3 - 2x_4 = -1\end{cases}$
    \item $\begin{cases} \sin(2x_1 + x_3) = \sqrt{5}\\ e^{-2x_2 - 2x_4} = 0\\ 4x_4 = 4\end{cases}$
  \end{enumerate}
  \begin{solucao}
    \begin{enumerate}[label={\alph*})]
      \item Não
    \item Sim, $x = \left(-\dfrac{5}{2 - 2\ln(2)} - 3\right)z$, $y = -\dfrac{5}{2 - 2\ln(2)}z$
      \item Sim, $x = 4$
      \item Sim, $x_1 = -\dfrac{5}{2} - \dfrac{7}{2}x_4$, $x_2 = -\dfrac{7}{10} - \dfrac{3}{5}x_3 - \dfrac{3}{10}x_4$
      \item Não
    \end{enumerate}
  \end{solucao}
\end{exercicio}

\begin{exercicio}
  Considere o sistema linear:
  \[
    \begin{cases}
      2x_1 - 4x_2 - x_3 = 1\\
      x_1 - 3x_2 + x_3 = 1\\
      3x_1 - 5x_2 - 3x_3 = 1
    \end{cases}
  \]
  Verifique se as seguintes ternas ordenadas são soluções desse sistema.
  \begin{enumerate}[label={\alph*})]
    \item $(3,1,1)$
    \item $(3,-1,1)$
    \item $(13, 5, 2)$
    \item $\left(\dfrac{13}{2}, \dfrac{5}{2}, 2\right)$
    \item $(17,7,5)$
  \end{enumerate}
  \begin{solucao}
    \begin{enumerate}[label={\alph*})]
      \item Sim
      \item Não
      \item Não
      \item Sim
      \item Sim
    \end{enumerate}
  \end{solucao}
\end{exercicio}

\begin{exercicio}
  Encontre a matriz ampliada dos seguintes sistemas lineares:
  \begin{enumerate}[label={\alph*})]
    \item $\begin{cases}x - y = 0\\2x + y = 3\end{cases}$
    \item $\begin{cases}2x_1 + 3x_2 - x_3 = 1\\x_1 + x_3 = 0\\-x_1 + 2x_2 - 2x_3 = 0\end{cases}$
    \item $\begin{cases}x + 5y = -1\\-x + y = -5\\2x + 4y = 4\end{cases}$
    \item $\begin{cases}a - 2b + d = 2\\-a + b - c - 3d = 1\end{cases}$
  \end{enumerate}
\end{exercicio}

\begin{exercicio}
  Nos itens abaixo encontre um sistema linear que tenha a matriz dada como sua matriz ampliada.
  \begin{multicols}{2}
    \begin{enumerate}[label={\alph*})]
      \item $\begin{amatrix}{3}0 & 1 & 1 & 1\\1 & -1 & 0 & 1\\2 & -1 & 1 & 1\end{amatrix}$
      \item $\begin{amatrix}{5}1 & -1 & 0 & 3 & 1 & 2\\1 & 1 & 2 & 1 & -1 & 4\\0 & 1 & 0 & 2 & 3 & 0\end{amatrix}$
      \item $\begin{pmatrix}7 & 2 & 1 & -3 & 5\\1 & 2 & 4 & \phantom{-} 0 & 1\end{pmatrix}$
      \item $\begin{pmatrix}\phantom{-} 2 & -1 \\ -4 & -6\\ \phantom{-} 1 & -1\\\phantom{-} 3 & \phantom{-} 0\end{pmatrix}$
      \item $\begin{pmatrix}2 & \phantom{-} 0 & 0\\3 & -4 &0\\0 & \phantom{-} 1 & 1\end{pmatrix}$
      \item $\begin{amatrix}{3}3 & 0 & -2 & 5\\7 & 1 & 4 & -3\\0 & -2 & 1 & 7\end{amatrix}$
      \item $\begin{amatrix}{3}1 & 2 & 3 & 4\\-4 & -3 & -2 & -1\\5 & -6 & 1 & 1\\-8 & 0 & 0 & 3\end{amatrix}$
      \item $\begin{pmatrix}\phantom{-} 3 & 0 & 1 & -4 & \phantom{-} 3\\-4 & 0 & 4 & \phantom{-} 1 & \phantom{-} 3\\-1 & 3 & 0 & -2 & -9\\\phantom{-} 0 & 0 & \phantom{-} 0 & -1 & -2\end{pmatrix}$
    \end{enumerate}
  \end{multicols}
\end{exercicio}

Nos Exercícios \ref{inicioverificamatrizes} à \ref{fimverificamatrizes}, determine se a matriz dada está: na forma escalonada, na forma escalonada reduzida por linhas, em ambas ou em nenhuma.

\begin{multicols}{2}
    \begin{exercicio}\label{inicioverificamatrizes}
        $\begin{bmatrix}1 & 0 & 0\\0 & 1 & 0\\0 & 0 & 1\end{bmatrix}$
        \begin{solucao}
            Ambas.
        \end{solucao}
    \end{exercicio}

        \begin{exercicio}
        $\begin{bmatrix}0 & 1 & 0\\1 & 0 & 0\\0 & 0 & 1\end{bmatrix}$
        \begin{solucao}
            Nenhuma.
        \end{solucao}
    \end{exercicio}

    \begin{exercicio}
        $\begin{bmatrix}1 & 0 & 0\\0 & 1 & 0\\0 & 0 & 0\end{bmatrix}$
        \begin{solucao}
            Ambas.
        \end{solucao}
    \end{exercicio}

    \begin{exercicio}
        $\begin{bmatrix}7 & 0 & \phantom{-} 1 & 4\\0 & 1 & -1 & 4\\0 & 0 & \phantom{-} 0 & 0\end{bmatrix}$
        \begin{solucao}
            Nenhuma.
        \end{solucao}
    \end{exercicio}

        \begin{exercicio}
        $\begin{bmatrix}1 & 2 & \phantom{-} 1 & 4\\0 & 1 & -1 & 4\\0 & 0 & \phantom{-} 0 & 0\end{bmatrix}$
        \begin{solucao}
            Forma escalonada.
        \end{solucao}
    \end{exercicio}

    \begin{exercicio}
        $\begin{bmatrix}0 & 1 & 0\\0 & 0 & 1\\0 & 0 & 0\end{bmatrix}$
        \begin{solucao}
            Ambas.
        \end{solucao}
    \end{exercicio}

    \begin{exercicio}
        $\begin{bmatrix}1 & 0 & 3 & 1\\0 & 1 & 2 & 4\end{bmatrix}$
        \begin{solucao}
            Ambas.
        \end{solucao}
    \end{exercicio}

    \begin{exercicio}
        $\begin{bmatrix}1 & 2 & 3\\1 & 0 & 0\\0 & 1 & 1\\0 & 0 & 1\end{bmatrix}$
        \begin{solucao}
            Forma escalonada.
        \end{solucao}
    \end{exercicio}

    \begin{exercicio}
    $\begin{bmatrix}1 & 1 & \phantom{-} 3 & \phantom{-} 5\\0 & 0 & -1 & -1\\0 & 0 & \phantom{-} 0 & \phantom{-} 1\end{bmatrix}$
    \begin{solucao}
        Nenhuma.
    \end{solucao}
\end{exercicio}

    \begin{exercicio}\label{fimverificamatrizes}
        $\begin{bmatrix}2 & 1 & 3 & \phantom{-} 5\\0 & 0 & 1 & -1\\0 & 0 & 0 & \phantom{-} 0\\0 & 0 & 0 & \phantom{-} 1\end{bmatrix}$
        \begin{solucao}
            Nenhuma.
        \end{solucao}
    \end{exercicio}
\end{multicols}


\begin{exercicio}
  Descrever explicitamente todas as possíveis formas que uma matriz $2 \times 2$ na forma escalonada reduzida por linhas pode assumir.
\end{exercicio}

\begin{exercicio}
  Descrever explicitamente todas as possíveis formas que uma matriz $3 \times 3$ na forma escalonada reduzida por linhas pode assumir.
\end{exercicio}

\begin{exercicio}
  Descrever explicitamente todas as possíveis formas que uma matriz $4 \times 4$ na forma escalonada reduzida por linhas pode assumir.
\end{exercicio}

\begin{exercicio}
  Nos itens a seguir mostre que as matrizes dadas são linha-equivalentes e encontre uma sequência de operações elementares com as
  linhas que convertem $A$ em $B$.
  \begin{enumerate}[label={\alph*})]
    \item $A = \begin{bmatrix}1 & 2\\3 & 4\end{bmatrix}$, $B = \begin{bmatrix}3 & 1\\1 & 0\end{bmatrix}$

    \item $A = \begin{bmatrix}\phantom{-} 2 & 0 & -1\\\phantom{-} 1 & 1 & \phantom{-} 0\\-1 & 1 & \phantom{-} 1\end{bmatrix}$,
      $B = \begin{bmatrix}3 & 1 & -1\\3 & 5 & \phantom{-} 1\\2 & 2 & \phantom{-} 0\end{bmatrix}$

      \item $A = \begin{bmatrix}\phantom{-} 2 & -3 & 10\\\phantom{-} 1 & -2 & 3\\-1 & \phantom{-} 3 & 1\end{bmatrix}$, $B = \begin{bmatrix}1 & -2 & 3\\0 & \phantom{-} 1 & 4\\0 & \phantom{-} 0 & 0\end{bmatrix}$

      \item $A = \begin{bmatrix}1 & 0 & 0 & \phantom{-} 2\\0 & 1 & 0 & \phantom{-} 3\\0 & 0 & 1 & -1\end{bmatrix}$, $B = \begin{bmatrix}\phantom{-} 1 & -2 & -3 & -1\\\phantom{-} 1 & -1 & -2 & \phantom{-} 1\\-1 & \phantom{-} 3 & \phantom{-} 5 & \phantom{-} 2\end{bmatrix}$
  \end{enumerate}
\end{exercicio}

Nos exercícios \ref{sistema_linear_inicio_eliminacao} à \ref{sistema_linear_eliminacao}, encontre a solução geral dos seguintes sistemas lineares usando o método de eliminação gaussiana.
\begin{exercicio}\label{sistema_linear_inicio_eliminacao}
  $
    \begin{cases}
      x + y + z = 4\\
      2x + 5y - 2z = 3\\
      x + 7y - 7z = 5
    \end{cases}
  $
  em $\rac$.
  \begin{solucao}
    Não existe solução.
  \end{solucao}
\end{exercicio}

\begin{exercicio}
  $
    \begin{cases}
      x - 2y + 3z = 0\\
      2x + 5y + 6z = 0
    \end{cases}
  $
  em $\real$.
  \begin{solucao}
    $S = \{(x, y, z) \mid z, y, z \in \real\} = \{(-3z, 0, z) \mid z \in \real\}$
  \end{solucao}
\end{exercicio}

\begin{exercicio}
  $
    \begin{cases}
      x_1 + x_2 + 2x_3 = 8\\
      -x_1 - 2x_2 + 3x_3 = 1\\
      3x_1 - 7x_2 + 4x_3 = 10
    \end{cases}
  $
  em $\real$.
  \begin{solucao}
    $S = \{(3, 1, 2)\}$
  \end{solucao}
\end{exercicio}

\begin{exercicio}
  $
    \begin{cases}
      2x_1 + 2x_2 + 2x_3 = 0\\
      -2x_1 + 5x_2 + 2x_3 = 1\\
      8x_1 + x_2 + 4x_3 = -1
    \end{cases}
  $
  em $\real$.
  \begin{solucao}
    $S = \{(x_1, x_2, x_3) \mid x_1, x_2, x_3 \in \real\} = \{(-1/7 - (3/7)x_3, 1/7 - (4/7)x_3, x_3) \mid x_3 \in \real\}$
  \end{solucao}
\end{exercicio}

\begin{exercicio}
  $
    \begin{cases}
      \phantom{2x_1} - 2x_2 + 3x_3 = 1\\
      3x_1 + 6x_2 - 3x_3 = -2\\
      6x_1 + 6x_2 + 3x_3 = 5
    \end{cases}
  $
  em $\real$.
  \begin{solucao}
    O sistema não tem solução.
  \end{solucao}
\end{exercicio}

\begin{exercicio}\label{sistema_linear_eliminacao}
  $
    \begin{cases}
      ix + iy = 0\\
      2ix - y = 0\\
    \end{cases}
  $
  em $\complex$.
  \begin{solucao}
    $S = \{(0, 0)\}$
  \end{solucao}
\end{exercicio}

Nos exercícios \ref{sistema_linear_inicio_gauss-jordan} à \ref{sistema_linear_fim_gauss-jordan}, encontre a solução geral dos seguintes sistemas lineares usando o método de eliminação de Gauss-Jordan.

\begin{exercicio}\label{sistema_linear_inicio_gauss-jordan}
  $
    \begin{cases}
      x_1 + x_2 + x_3 + x_4 = 0\\
      x_1 + x_2 + x_3 - x_4 = 4\\
      x_1 + x_2 - x_3 + x_4 = -4\\
      x_1 - x_2 + x_3 + x_4 = 2\\
    \end{cases}
  $
  em $\real$.
  \begin{solucao}
    $S = \{(1, -1, 2, -2)\}$
  \end{solucao}
\end{exercicio}

\begin{exercicio}
  $
    \begin{cases}
      -2x_2 + 3x_3 = 1\\
      3x_1 + 6x_2 - 3x_3 = -1\\
      6x_1 + 6x_2 + 3x_3 = 5
    \end{cases}
  $
  em $\real$.
  \begin{solucao}
    Não existe solução.
  \end{solucao}
\end{exercicio}

\begin{exercicio}
  $
    \begin{cases}
      2x_1 + x_2 + (6 + 6i)x_3 + 8x_4 = 0\\
      x_1 + x_2 + (2 + 5i)x_3 + (5  i)x_4 = 0\\
      (2 + 2i)x_1 + 2x_2 + (2 + 14i)x_3 + (14 + 8i)x_4 =0\\
      (-1 - i)x_1 - x_2 - (3 + 3i)x_3 + (-6 + 2i)x_4 = 0
    \end{cases}
  $
  em $\complex$.
  \begin{solucao}
    $S = \{(x_1, x_2, x_3, x_4) \mid x_1, x_2, x_3, x_4 \in \complex\} = $\\ $\{(-(2 + i)x_3 - (3 - i)x_4,-3ix_3 - 2x_4, x_3, x_4) \mid x_3, x_4 \in \complex\}$
  \end{solucao}
\end{exercicio}

\begin{exercicio}
  $
    \begin{cases}
      x_1 + 2x_2 - 3x_4 + x_5 = 2\\
      x_1 + 2x_2 + x_3 - 3x_4 + x_5 + 2x_6 = 3\\
      x_1 + 2x_2 - 3x_4 + 2x_5 + x_6 = 4\\
      3x_1 + 6x_2 + x_3 - 9x_4 + 4x_5 + 3x_6 = 9
    \end{cases}
  $
  em $\real$.
  \begin{solucao}
    $S = \{(x_1, x_2, x_3, x_4, x_5, x_6) \mid x_1, x_2, x_3, x_4, x_5, x_6 \in \real\} = $\\ $\{(x_6 + 3x_4 - 2x_2, x_2, 1 - 2x_6, x_4, 2 - x_6, x_6) \mid x_2, x_4, x_6 \in \real\}$
  \end{solucao}
\end{exercicio}

\begin{exercicio}
  $
    \begin{cases}
      x_1 + 3x_2 - 2x_3 + 2x_5 = 0\\
      2x_1 + 6x_2 - 5x_3 - 2x_4 + 4x_5 - 3x_6 = -1\\
      5x_3 + 10x_4 + 15x_6 = 5\\
      2x_1 + 6x_2 + 8x_4 + 4x_5 + 18x_6 = 6
    \end{cases}
  $
  em $\real$.
  \begin{solucao}
    $S = \{(x_1, x_2, x_3, x_4, x_5, x_6) \mid x_1, x_2, x_3, x_4, x_5, x_6 \in \real\} = $\\ $\{(-3x_2 - 4x_4 - 2x_5, x_2, -2x_4, x_4, x_5, 1/3) \mid x_2, x_4, x_5 \in \real\}$
  \end{solucao}
\end{exercicio}

\begin{exercicio}\label{sistema_linear_fim_gauss-jordan}
  $
    \begin{cases}
      3x + 2y - 4z = 1\\
      x - y + z = 3\\
      x - y - 3z = -3\\
      3x + 3y - 5z =0\\
      -x + y + z = 1
    \end{cases}
  $
  em $\rac$.
\begin{solucao}
  Não existe solução.
\end{solucao}
\end{exercicio}

\begin{exercicio}
  Determine o número de soluções do sistema
  \[
    \begin{cases}
      x + 2y - 3z = 4\\
      4x + y + 2z = 6\\
      x + 2y + (\alpha^2 - 19)z = \alpha,
    \end{cases}
\]
em função do parâmetro:
\begin{enumerate}[label={\alph*})]
   \item $\alpha \in \real$;
   \item $\alpha \in \rac$?
 \end{enumerate}
\end{exercicio}

\begin{exercicio}
  Considere o sistema linear
  \[
    \begin{cases}
      (k + 2)x + 2ky - z = 1\\
      x - 2y + kz = -k\\
      y + z = k
    \end{cases}.
  \]
  \begin{enumerate}[label={\alph*})]
    \item Para quais valores de $k \in \real$ o sistema admite solução única? Determine a solução, caso exista.
    \item Para quais valores de $k \in \real$ o sistema admite infinitas soluções? Determine tais soluções, caso existam.
  \end{enumerate}
\end{exercicio}

Nos exercícios \ref{sistemasinicio} à \ref{sistemasfim}, encontre o(s) valor(es) de $\lambda$ tal(is) que:
\begin{enumerate}[label={\alph*})]
  \item o sistema tem solução única;
  \item o sistema tem várias soluções;
  \item o sistema não tem solução.
\end{enumerate}

\begin{exercicio}\label{sistemasinicio}
  $
    \begin{cases}
      x + \lambda y = 1\\
      \lambda x + y = 1
    \end{cases}
  $
  em $\real$.
  \begin{solucao}
    \begin{enumerate}[label={\alph*})]
      \item $\lambda \ne -1$ e $\lambda \ne 1$
      \item $\lambda = 1$
      \item $\lambda = -1$
    \end{enumerate}
  \end{solucao}
\end{exercicio}

\begin{exercicio}
  $
    \begin{cases}
      (\lambda - 2)x + y = 0\\
      x + (\lambda - 2)y = 0
    \end{cases}
  $
  em $\real$.
  \begin{solucao}
    \begin{enumerate}[label={\alph*})]
      \item $\lambda \ne 1$ e $\lambda \ne 3$
      \item $\lambda = 1$ e $\lambda = 3$
      \item Não existe tal $\lambda$
    \end{enumerate}
  \end{solucao}
\end{exercicio}

\begin{exercicio}
  $
    \begin{cases}
      x + 2y - 3z = 4\\
      3x - y + 5z = 2\\
      4x + y + (\lambda^2 - 14)z = \lambda + 2
    \end{cases}
  $
  em $\real$.
  \begin{solucao}
    \begin{enumerate}[label={\alph*})]
      \item $\lambda \ne \pm 4$
      \item $\lambda = 4$
      \item $\lambda = -4$
    \end{enumerate}
  \end{solucao}
\end{exercicio}

\begin{exercicio}
  $
    \begin{cases}
      x + y + z = 2\\
      2x + 3y + 2z = 5\\
      2x + 3y + (\lambda^2 - 1)z = \lambda + 1
    \end{cases}
  $
  em $\rac$.
  \begin{solucao}
    \begin{enumerate}[label={\alph*})]
      \item Não existe tal $\lambda$.
      \item Não existe tal $\lambda$.
      \item Não existe tal $\lambda$.
    \end{enumerate}
  \end{solucao}
\end{exercicio}

\begin{exercicio}
  $
    \begin{cases}
      x_1 + \lambda x_2 + (1 + 4\lambda )x_3 = 1 + 4\lambda \\
      2x_1 + (\lambda  + 1)x_2 + (2 + 7\lambda )x_3 = 1 + 7\lambda \\
      3x_1 + (\lambda  + 2)x_2 + (3 + 9\lambda )x_3 = 1 + 9\lambda
    \end{cases}
  $
  \begin{solucao}
    \begin{enumerate}[label={\alph*})]
      \item $\lambda \ne 0$ e $\lambda \ne 1$
      \item $\lambda = 0$
      \item $\lambda = 1$
    \end{enumerate}
  \end{solucao}
\end{exercicio}

\begin{exercicio}\label{sistemasfim}
  $
    \begin{cases}
      x + y = 2\\
      y + z = 2\\
      x + z = 2\\
      x + y + \lambda z = 0
    \end{cases}
  $
  em $\real$.
  \begin{solucao}
    \begin{enumerate}[label={\alph*})]
      \item $\lambda = -2$
      \item Não existe tal $\lambda$.
      \item $\lambda \ne -2$
    \end{enumerate}
  \end{solucao}
\end{exercicio}

\begin{exercicio}
  Encontre condições sobre os $b_i \in \real$ para que cada um dos sistemas tenha solução.
  \begin{enumerate}[label={\alph*})]
    \item $\begin{cases}
      x_1 - 2x_2 + 5x_3 = b_1\\
      4x_1 - 5x_2 + 8x_3 = b_2\\
      -3x_1 + 3x_2 - 3x_3 = b_3
    \end{cases}$

    \item $\begin{cases}
      x_1 - 2x_2 - x_3 = b_1\\
      -4x_1 + 5x_2 + 2x_3 = b_2\\
      -4x_1 + 7x_2 + 4x_3 = b_3
    \end{cases}$
  \end{enumerate}
  \begin{solucao}
    \begin{enumerate}[label={\alph*})]
      \item $b_3 - b_1 + b_2 = 0$
      \item O sistema tem solução para todos os valores reais de $b_1$, $b_2$ e $b_3$.
    \end{enumerate}
  \end{solucao}
\end{exercicio}

\newpage
\Closesolutionfile{ans}
\hrule
\begin{center}
{\large\bf RESPOSTAS}
\end{center}
\hrule
\input{ans1}

\end{document}